%! Principles of Experimental Design — Unit 1, Lesson 1.6 — Guided Notes (Answer Key)
\documentclass[10pt]{article}
\usepackage{saar-article}
\usepackage{saar-key}   % pulls in -boxes; do NOT also load -boxes
\newcommand{\TallMath}[1]{$\displaystyle #1\rule[-1.4em]{0pt}{3.2em}$}

\begin{document}

\pageheader{Unit 1, Lesson 1.6}{Guided Notes --- Answer Key}
% NAMESTRIP: no \namedateperiod here — mirrors the blank. Teacher-only prose goes
% in the lesson plan as \begin{teachernote}[Guided Notes], never in a key.

\vspace{0.15in}

\begin{objectivebox}
\small
By the end of this lesson, I will be able to\dots
\begin{enumerate}[leftmargin=1.5em, itemsep=2pt, topsep=2pt]
  \item explain what makes a study an \textbf{experiment} --- the researcher
        assigns the groups and imposes a treatment;
  \item name the \textbf{experimental units}, the \textbf{treatment}, and the
        \textbf{control group} in a described experiment;
  \item explain \textbf{blinding} and the \textbf{placebo effect}, and say who
        is blinded in a design;
  \item tell apart a \textbf{completely randomized} design, a \textbf{block}
        design, and \textbf{matched pairs};
  \item evaluate a design against the four principles ---
        \textbf{comparison, randomization, replication, control} --- and justify
        which conclusions it supports.
\end{enumerate}
\end{objectivebox}

\boxguard
\begin{vocabbox}
\small
Fill in each term as we build it together.
\vocabans{Experiment}{A study in which the researcher assigns the experimental
units to groups and imposes a treatment on at least one group. Because the
researcher builds the groups, an experiment can support a cause-and-effect
conclusion.}
\vocabans{Experimental units (subjects)}{The individuals or objects a treatment
is applied to --- at Riverbend, the $120$ student volunteers. When the units are
people they are called subjects.}
\vocabans{Treatment}{The condition the researcher applies to a group --- at
Riverbend, attending the Study Center twice a week for the quarter.}
\vocabans{Control group}{The group that does not receive the treatment. It gives
a baseline showing what happens without the treatment, so there is something to
compare the treatment group against.}
\vocabans{Random assignment}{Using chance --- a generator, a coin, drawn slips
--- to decide which group each unit goes in, so the groups start out alike and
no one chooses their own group.}
\vocabans{Blinding}{Keeping someone in the experiment from knowing which group a
subject is in. Single-blind: the subjects (or the measurers) do not know.
Double-blind: neither does.}
\end{vocabbox}

\boxguard[12]
\begin{hookbox}
\small
Last class the Riverbend counselor \emph{watched}: of $60$ students who used the
Study Center, $65\%$ passed every class, against $80\%$ of the $90$ who never
went. The Study Center looked harmful. This quarter she \emph{assigned} instead
--- $120$ volunteers, split by a random number generator --- and got $75\%$
against $55\%$. The Study Center looks helpful.

\vspace{2pt}
\textbf{Same school, same Study Center, opposite conclusions. Which study should
she believe, and what makes it different?}
\ansline{Answers vary. The second one --- she built the groups by chance, so struggling}
\ansline{students were not piling into one of them. In the first, the students chose.}
\end{hookbox}

\boxguard[30]
\begin{notesbox}{1. Building the Groups Instead of Watching Them}
\small
In an \textbf{experiment} the researcher does not wait to see what happens. The
researcher \textbf{assigns} individuals to groups and then \textbf{imposes a
treatment} on at least one group.

\vspace{3pt}
Three words name the pieces. The individuals or objects a treatment is applied
to are the \textbf{experimental units}; when they are people we call them
\textbf{subjects}. What is applied to a group is the \textbf{treatment}. The
group that does \emph{not} get the treatment is the \textbf{control group}, and
it exists to give a baseline --- what happens \emph{without} the treatment.

\vspace{3pt}
Which of these plans are experiments?

\begin{center}
\renewcommand{\arraystretch}{1.4}
\begin{tabularx}{0.98\linewidth}{@{} X c p{3.4cm} @{}}
\toprule
\textbf{The plan} & \textbf{Experiment?} & \textbf{Who decided the groups?} \\
\midrule
Pull $150$ student records and compare the pass rate of students who used the
Study Center to those who did not. & \ans{no} & \ans{the students did} \\
Ask for volunteers, then let a random number generator send $60$ to the Study
Center and $60$ nowhere. & \ans{yes} & \ans{chance, run by the counselor} \\
Survey a random sample of pool members about what time they swim and how well
they sleep. & \ans{no} & \ans{the members did} \\
Give $60$ randomly chosen volunteers a new tutoring app and the other $60$ the
old worksheet packet. & \ans{yes} & \ans{chance, run by the researcher} \\
\bottomrule
\end{tabularx}
\end{center}

\vspace{3pt}
In Riverbend's experiment, the experimental units are \ans{the $120$ volunteers},
the treatment is \ans{attending twice a week}, and the control group is
\ans{the $60$ told not to attend}.

\vspace{2pt}
\textit{Lesson 1.5 turned on one word --- \emph{assign}. So does this one. What
changed is which side of it you are standing on.}
\end{notesbox}

\boxguard[30]
\begin{notesbox}{2. The Four Principles of Experimental Design}
\small
Almost every experiment worth trusting does all four of these. Fill in the last
column for Riverbend's design.

\begin{center}
\renewcommand{\arraystretch}{1.4}
\begin{tabularx}{0.98\linewidth}{@{} p{2.4cm} X p{4.1cm} @{}}
\toprule
\textbf{Principle} & \textbf{What the design must do}
  & \textbf{How Riverbend does it} \\
\midrule
Comparison & Use at least two groups, one of them a control group, so there is
  something to compare against. & \ans{$60$ attend, $60$ do not} \\
Randomization & Let \emph{chance} --- not the researcher, not the subject ---
  decide which group each unit goes in. & \ans{a random generator splits the
  $120$} \\
Replication & Put enough units in each group that one person's good or bad
  luck cannot decide the result. & \ans{$60$ students per group, not $1$} \\
Control & Keep every \emph{other} condition the same for both groups.
  & \ans{same quarter, classes, teachers} \\
\bottomrule
\end{tabularx}
\end{center}

\vspace{3pt}
A variable that differs between the groups \emph{and} also affects the response
is a \textbf{confounding variable}. When one is present you cannot tell whether
the difference came from the treatment or from it. In Lesson 1.5 you called such
a variable a \ans{lurking} variable, because it sat outside the study; the
principle that keeps it out of an experiment is \ans{control}.

\vspace{2pt}
\textit{Randomization makes the two groups alike \emph{before} you start.
Control keeps them alike \emph{while} you run.}
\end{notesbox}

\boxguard[30]
\begin{notesbox}{3. Blinding and the Placebo Effect}
\small
Suppose the Study Center students know they were picked for a special program
and try harder for that reason alone. Then part of the $20$-point gap is not the
tutoring --- it is \emph{being chosen}. When a subject responds to the
\emph{idea} of a treatment rather than to the treatment itself, that is the
\textbf{placebo effect}. A fake treatment built to look real is a
\textbf{placebo}.

\vspace{3pt}
\textbf{Blinding} means somebody involved in the experiment is not told which
group a subject is in.

\begin{center}
\renewcommand{\arraystretch}{1.4}
\begin{tabularx}{0.98\linewidth}{@{} X c c @{}}
\toprule
\textbf{Version of the design} & \textbf{Students blinded?}
  & \textbf{Graders blinded?} \\
\midrule
Only the treatment group meets after school; the teachers who grade see the
attendance roster. & \ans{no} & \ans{no} \\
Only the treatment group meets after school; the teachers who grade are never
told who attended. & \ans{no} & \ans{yes} \\
\emph{Both} groups meet after school --- one for tutoring, one for a quiet
study hall --- and the graders are never told which room was which.
  & \ans{yes} & \ans{yes} \\
\bottomrule
\end{tabularx}
\end{center}

\vspace{3pt}
The quiet study hall in the last row is a \ans{placebo}: it keeps the control
group from knowing it is the control group. A design in which
\emph{neither} the subjects nor the people measuring the response know is called
\ans{double-blind}.

\vspace{2pt}
\textit{Blinding does not change who got the treatment. It removes the part of
the response that comes from \emph{knowing}.}
\end{notesbox}

\boxguard[30]
\begin{notesbox}{4. Three Ways to Assign the Units}
\small
\begin{enumerate}[leftmargin=1.6em, itemsep=3pt, topsep=3pt]
  \item \textbf{Completely randomized design} --- every unit is assigned to a
        treatment purely by chance. Riverbend's design is this one.
  \item \textbf{Randomized block design} --- first sort the units into
        \textbf{blocks} of \emph{similar} units, then randomize \emph{inside}
        each block.
  \item \textbf{Matched pairs} --- the block is a pair: two very similar units,
        one sent to each treatment. Sometimes a unit is paired with
        \emph{itself}, measured before and after.
\end{enumerate}

\vspace{3pt}
The counselor worries that strong and struggling students may not land evenly.
Of the $120$ volunteers, $80$ passed every class last quarter and $40$ did not,
so she blocks on that and randomizes inside each block.

\begin{center}
\renewcommand{\arraystretch}{1.4}
\begin{tabularx}{0.95\linewidth}{@{} X c c c @{}}
\toprule
\textbf{Block} & \textbf{In block} & \textbf{To Study Center}
  & \textbf{To control} \\
\midrule
Passed every class last quarter & $80$ & \ans{$40$} & \ans{$40$} \\
Did not pass every class last quarter & $40$ & \ans{$20$} & \ans{$20$} \\
\midrule
Total & $120$ & \ans{$60$} & \ans{$60$} \\
\bottomrule
\end{tabularx}
\end{center}

\vspace{3pt}
Blocking does not erase the difference between strong and struggling students.
What it does is \ans{guarantee each group gets the same mix of them}.

\vspace{2pt}
\textit{Block on the variable you are most afraid of. Randomize everything
else.}
\end{notesbox}

\boxguard[30]
\begin{notesbox}{5. What This Experiment Can Claim}
\small
The quarter ends. Here is the whole experiment in one table --- fill in the
percents from your warm-up.

\begin{center}
\renewcommand{\arraystretch}{1.4}
\begin{tabularx}{0.95\linewidth}{@{} X c c c @{}}
\toprule
\textbf{Group} & \textbf{Size} & \textbf{Passed every class} & \textbf{Percent} \\
\midrule
Study Center (treatment) & $60$ & $45$ & \ans{$75\%$} \\
No Study Center (control) & $60$ & $33$ & \ans{$55\%$} \\
\midrule
All volunteers & $120$ & $78$ & \ans{$65\%$} \\
\bottomrule
\end{tabularx}
\end{center}

\vspace{3pt}
The treatment group is higher by \ans{$20$} percentage points. Now compare
what each of Riverbend's two studies is allowed to say:

\begin{center}
\renewcommand{\arraystretch}{1.45}
\begin{tabularx}{0.98\linewidth}{@{} p{4.3cm} p{3.5cm} X @{}}
\toprule
\textbf{The study} & \textbf{What it found} & \textbf{What it may claim} \\
\midrule
Lesson 1.5 --- watched $150$ records & $65\%$ of users vs.\ $80\%$ of non-users
  & \ans{an association only --- the students chose} \\
Lesson 1.6 --- assigned $120$ volunteers & $75\%$ treatment vs.\ $55\%$ control
  & \ans{the Study Center caused the higher pass rate} \\
\bottomrule
\end{tabularx}
\end{center}

\vspace{3pt}
The reason fits in one sentence: because \ans{chance assigned the groups}, the
two groups were alike at the start, so the treatment is the only thing left that
can explain the gap.

\vspace{2pt}
\textit{Watching earns you an association. Assigning is what buys the word
\emph{caused}.}
\end{notesbox}

\boxguard
\begin{practicebox}
\small
\textbf{Cedar Ridge Apartments.} $100$ households volunteer for a study of
whether a weekly recycling reminder increases recycling. The manager randomly
assigns $50$ households to receive the reminder and $50$ to receive nothing,
then weighs each household's recycling for a month.

\vspace{3pt}
Experimental units: \ans{the $100$ households} \hfill Treatment:
\ans{the weekly reminder} \hfill Control: \ans{the $50$ with no reminder}

\vspace{3pt}
Each row changes the plan. Name the principle the change breaks.

\begin{center}
\renewcommand{\arraystretch}{1.45}
\begin{tabularx}{0.98\linewidth}{@{} X p{3.6cm} @{}}
\toprule
\textbf{Instead, the manager\dots} & \textbf{Principle broken} \\
\midrule
sends the reminder to all $100$ households and checks how much they recycle.
  & \ans{comparison} \\
sends it to the first $50$ households that signed up. & \ans{randomization} \\
sends it to $1$ household and nothing to $1$ other household. & \ans{replication} \\
sends it to the $50$ households in the renovated building, which just got new
recycling bins, and nothing to the $50$ in the old building. & \ans{control} \\
\bottomrule
\end{tabularx}
\end{center}

\vspace{3pt}
With the original plan, the reminder group recycles more. May the manager write
``the reminder increased recycling''? Say why or why not.
\ansline{Yes. Chance decided which households got the reminder, so the two groups started}
\ansline{out alike and everything else was kept the same --- the reminder is what is left.}
\end{practicebox}

\end{document}
